\ifdefined\gkver\else\def\gkver{student}\fi
\documentclass[\gkver]{gaokaozhenti}
\begin{document}

\chapter{2011年江苏}

\begin{enumerate}
\item
%% number: 1
%% paperName: 2011年江苏高考第 1 题 3 分
%% typeId: 1
%% type: 单选题
%% chapter: 力物体的平衡
%% point: 三力平衡
%% method: 无
%% score: 3
%% degree: 650
%% duplicateId: 0
%% body:
如图所示，石拱桥的正中央有一质量为 $m$ 的对称楔形石块，侧面与竖直方向的夹角为 $\alpha$，重力加速度为 $g$，若接触面间的摩擦力忽略不计，则石块侧面所受弹力的大小为\xzanswer{A}

\onepicture[w=6.0cm]{figs/01.png}

\fourchoices[answer=A]
{$\frac{mg}{2\sin\alpha}$}
{$\frac{mg}{2\cos\alpha}$}
{$\frac{1}{2}mg\tan\alpha$}
{$\frac{1}{2}mg\cot\alpha$}

%% answer: A
%% memo:
\memoanswer{
楔形石块受力如图，根据力的合成可得：

$mg=2\times F\cos(90^{\circ}-\alpha)$，所以 $F=\frac{mg}{2\cos(90^{\circ}-\alpha)}=\frac{mg}{2\sin\alpha}$，A 正确。

\onepicture[w=4.6cm]{figs/01_an.png}
}

\item
%% number: 2
%% paperName: 2011年江苏高考第 2 题 3 分
%% typeId: 1
%% type: 单选题
%% chapter: 电磁感应
%% point: 切割磁感线电动势
%% method: 无
%% score: 3
%% degree: 650
%% duplicateId: 0
%% body:
如图所示，固定的水平长直导线中通有电流 $I$，矩形线框与导线在同一竖直平面内，且一边与导线平行。线框由静止释放，在下落过程中\xzanswer{B}

\onepicture[w=4.4cm]{figs/02.png}

\fourchoices[answer=B]
{穿过线框的磁通量保持不变}
{线框中感应电流方向保持不变}
{线框所受安培力的合力为零}
{线框的机械能不断增大}

%% answer: B
%% memo:
\memoanswer{
因为磁感应强度随线框下落而减小，所以磁通量也减小，A 错误；因为磁通量随线框下落而减小，根据楞次定律，感应电流的磁场与原磁场方向相同，不变，所以感应电流的方向不变，本题选 B；感应电流在磁场中受安培力作用，上框边比下框边始终处于较强的磁场区域，线框所受安培力的合力向上不为零，C 错误；下落过程中克服安培力做功，机械能转化为内能，机械能减少，D 错误。
}

\item
%% number: 3
%% paperName: 2011年江苏高考第 3 题 3 分
%% typeId: 1
%% type: 单选题
%% chapter: 直线运动
%% point: 运动的合成
%% method: 无
%% score: 3
%% degree: 450
%% duplicateId: 0
%% body:
如图所示，甲、乙两同学从河中 O 点出发，分别沿直线游到 A 点和 B 点后，立即沿原路线返回到 O 点，OA、OB 分别与水流方向平行和垂直，且 $OA=OB$。若水流速度不变，两人在静水中游速相等，则他们所用时间 $t_{\text{甲}}$、$t_{\text{乙}}$ 的大小关系为\xzanswer{C}

\onepicture[w=5.0cm]{figs/03.png}

\fourchoices[answer=C]
{$t_{\text{甲}}<t_{\text{乙}}$}
{$t_{\text{甲}}=t_{\text{乙}}$}
{$t_{\text{甲}}>t_{\text{乙}}$}
{无法确定}

%% answer: C
%% memo:
\memoanswer{
设游速为 $v$，水速为 $v_{0}$，$OA=OB=l$，则甲时间 $t_{\text{甲}}=\frac{l}{v+v_{0}}+\frac{l}{v-v_{0}}$；

乙沿 OB 运动，乙的速度矢量图如图，合速度必须沿 OB 方向，则乙时间 $t_{\text{乙}}=\frac{l}{\sqrt{v^{2}-v_{0}^{2}}}\times2$，联立解得：$t_{\text{甲}}>t_{\text{乙}}$，C 正确。

\onepicture[w=3.8cm]{figs/03_an.png}
}

\item
%% number: 4
%% paperName: 2011年江苏高考第 4 题 3 分
%% typeId: 1
%% type: 单选题
%% chapter: 功和能
%% point: 动能定理
%% method: 近似与估算
%% score: 3
%% degree: 650
%% duplicateId: 0
%% body:
如图所示，演员正在进行杂技表演。由图可估算出他将一只鸡蛋抛出的过程中对鸡蛋所做的功最接近于\xzanswer{A}

\onepicture[w=5.0cm]{figs/04.png}

\fourchoices[answer=A]
{$0.3\UJ$}
{$3\UJ$}
{$30\UJ$}
{$300\UJ$}

%% answer: A
%% memo:
\memoanswer{
生活经验告诉我们：10 个鸡蛋大约 1 斤即 $0.5\Ukg$，则一个鸡蛋的质量约为 $m=\frac{0.5}{10}=0.05\Ukg$，鸡蛋大约能抛高度 $h=0.6\Um$，则做功约为 $W=mgh=0.05\times10\times0.6\UJ=0.3\UJ$，A 正确。
}

\item
%% number: 5
%% paperName: 2011年江苏高考第 5 题 3 分
%% typeId: 1
%% type: 单选题
%% chapter: 电磁感应
%% point: 电磁感应中的变加速
%% method: 无
%% score: 3
%% degree: 250
%% duplicateId: 0
%% body:
如图所示，水平面内有一平行金属导轨，导轨光滑且电阻不计。匀强磁场与导轨平面垂直。阻值为 $R$ 的导体棒垂直于导轨静止放置，且与导轨接触。$t=0$ 时，将开关 S 由 1 掷到 2。$q$、$i$、$v$ 和 $a$ 分别表示电容器所带的电荷量、棒中的电流、棒的速度和加速度。下列图象正确的是\xzanswer{D}

\onepicture[w=5.2cm]{figs/05.png}

\fourchoices[answer=D, ispicture=true, h=2.4cm]
{figs/05a.png}
{figs/05b.png}
{figs/05c.png}
{figs/05d.png}

%% answer: D
%% memo:
\memoanswer{
$t=0$ 时，将开关 S 由 1 掷到 2，电容器放电，开始时 $i=\frac{E}{R}$，因安培力作用使导体棒产生加速度，导体棒速度增大，产生反向感应电动势，使电流减小，安培力减小，加速度减小，减小至零时，速度达最大值 $v_{m}$ 做匀速运动，电容器两极电压为 $BLv_{m}$（$L$ 为导轨宽度），A、B、C 错误，D 正确。
}

\item
%% number: 6
%% paperName: 2011年江苏高考第 6 题 4 分
%% typeId: 2
%% type: 多选题
%% chapter: 电路
%% point: 传感器
%% method: 无
%% score: 4
%% degree: 650
%% duplicateId: 0
%% body:
美国科学家 Willard S. Boyle 与 George E. Smith 因电荷耦合器件（CCD）的重要发明荣获 2009 年度诺贝尔物理学奖。CCD 是将光学量转变成电学量的传感器。下列器件可作为传感器的有\xzanswer{BC}

\fourchoices[answer=BC]
{发光二极管}
{热敏电阻}
{霍尔元件}
{干电池}

%% answer: BC
%% memo:
\memoanswer{
热敏电阻的阻值随温度的变化而变化，可以把温度转化为电学量，而霍尔元件可以把磁场的磁感应强度转化为电学量，热敏电阻和霍尔元件可作为传感器，B、C 正确。
}

\item
%% number: 7
%% paperName: 2011年江苏高考第 7 题 4 分
%% typeId: 2
%% type: 多选题
%% chapter: 曲线运动
%% point: 天体运动
%% method: 无
%% score: 4
%% degree: 250
%% duplicateId: 0
%% body:
一行星绕恒星作圆周运动。由天文观测可得，其运动周期为 $T$，速度为 $v$，引力常量为 $G$，则\xzanswer{ACD}

\fourchoices[answer=ACD]
{恒星的质量为 $\frac{v^{3}T}{2\pi G}$}
{行星的质量为 $\frac{4\pi^{2}v^{3}}{GT^{2}}$}
{行星运动的轨道半径为 $\frac{vT}{2\pi}$}
{行星运动的加速度为 $\frac{2\pi v}{T}$}

%% answer: ACD
%% memo:
\memoanswer{
根据 $F=G\frac{Mm}{r^{2}}=m\left(\frac{2\pi}{T}\right)^{2}$、$v=\frac{2\pi r}{T}$ 得：$M=\frac{v^{3}T}{2\pi G}$、$r=\frac{vT}{2\pi}$，A、C 正确，B 错误；根据 $a=\frac{v^{2}}{r}$、$v=\omega r=\frac{2\pi}{T}r$ 得：$a=\frac{2\pi v}{T}$，D 正确。
}

\item
%% number: 8
%% paperName: 2011年江苏高考第 8 题 4 分
%% typeId: 2
%% type: 多选题
%% chapter: 电场
%% point: 带电粒子的偏转
%% method: 无
%% score: 4
%% degree: 450
%% duplicateId: 0
%% body:
一粒子从 A 点射入电场，从 B 点射出，电场的等势面和粒子的运动轨迹如图所示，图中左侧前三个等势面彼此平行，不计粒子的重力。下列说法正确的有\xzanswer{AB}

\onepicture[w=5.4cm]{figs/08.png}

\fourchoices[answer=AB]
{粒子带负电荷}
{粒子的加速度先不变，后变小}
{粒子的速度不断增大}
{粒子的电势能先减小，后增大}

%% answer: AB
%% memo:
\memoanswer{
由于电场线与等势面垂直（如图），电场线先向右后向上偏，而粒子却向下偏了，所以电场力与电场强度方向相反，所以粒子带负电，A 正确；又等势面先平行并且密集，后变稀疏，说明电场强度先不变，后变小，则粒子受电场力先不变，后变小，所以加速度先不变，后变小，B 正确；电场力与初速度方向相反，所以速度先减小，C 错误；而电场力先做负功，所以电势能先增大，D 错误。

\onepicture[w=5.2cm]{figs/08_an.png}
}

\item
%% number: 9
%% paperName: 2011年江苏高考第 9 题 4 分
%% typeId: 2
%% type: 多选题
%% chapter: 运动和力
%% point: 连接体
%% method: 无
%% score: 4
%% degree: 250
%% duplicateId: 0
%% body:
如图所示，倾角为 $\alpha$ 的等腰三角形斜面固定在水平面上，一足够长的轻质绸带跨过斜面的顶端铺放在斜面的两侧，绸带与斜面间无摩擦。现将质量分别为 $M$、$m$（$M>m$）的小物块同时轻放在斜面两侧的绸带上。两物块与绸带间的动摩擦因数相等，且最大静摩擦力与滑动摩擦力大小相等。在 $\alpha$ 角取不同值的情况下，下列说法正确的有\xzanswer{AC}

\onepicture[w=6.0cm]{figs/09.png}

\fourchoices[answer=AC]
{两物块所受摩擦力的大小总是相等}
{两物块不可能同时相对绸带静止}
{$M$ 不可能相对绸带发生滑动}
{$m$ 不可能相对斜面向上滑动}

%% answer: AC
%% memo:
\memoanswer{
由于绸带与斜面之间光滑，并且 $M>m$，所以 M、m 和绸带一起向左滑动，加速度为 $a$，由牛顿第二定律对整体：$Mg\sin\alpha-mg\sin\alpha=(M+m)a$；对 M 有：$Mg\sin\alpha-F_{fM}=Ma$；对 m 有：$F_{fm}-mg\sin\alpha=ma$，解得 $F_{fM}=F_{fm}=\frac{2Mm}{M+m}g\sin\alpha$，即 A、C 正确，B、D 错误。
}

\item
%% number: 10
%% paperName: 2011年江苏高考第 10 题 8 分
%% typeId: 4
%% type: 实验题
%% chapter: 力物体的平衡
%% point: 验证平行四边形实验
%% method: 无
%% score: 8
%% degree: 650
%% duplicateId: 0
%% body:
某同学用如图所示的实验装置来验证“力的平行四边形定则”。弹簧测力计 A 挂于固定点 P，下端用细线挂一重物 $M$。弹簧测力计 B 的一端用细线系于 O 点，手持另一端向左拉，使结点 O 静止在某位置。分别读出弹簧测力计 A 和 B 的示数，并在贴于竖直木板的白纸上记录 O 点的位置和拉线的方向。

\begin{enumerate}
	\item
	本实验用的弹簧测力计示数的单位为 $\UN$，图中 A 的示数为\tkanswer[1.5cm]{3.6} $\UN$。
	\item
	下列不必要的实验要求是\tkanswer[1.0cm]{D}。（请填写选项前对应的字母）

	（A）应测量重物 $M$ 所受的重力

	（B）弹簧测力计应在使用前校零

	（C）拉线方向应与木板平面平行

	（D）改变拉力，进行多次实验，每次都要使 O 点静止在同一位置
	\item
	某次实验中，该同学发现弹簧测力计 A 的指针稍稍超出量程，请你提出两个解决办法。
\end{enumerate}

\onepicture[w=6.4cm, align=right]{figs/10.png}

%% answer:
\jdanswer{
\begin{enumerate}
	\item
	3.6
	\item
	D
	\item
	①改变弹簧测力计 B 拉力的大小；②减小重物 $M$ 的质量（或将 A 更换为较大的测力计，改变弹簧测力计 B 拉力的方向）
\end{enumerate}
}
%% memo:
\memoanswer{
（1）如图，弹簧测力计 A 的示数为 $3.6\UN$；

（2）“验证平行四边形定则”，只要验证每次 $F_{\mathrm{A}}$、$F_{\mathrm{B}}$ 和 $Mg$ 满足平行四边形即可，D 不必要。

（3）弹簧测力计 A 的指针稍稍超出量程，说明拉弹簧测力计 A 的力大了，由力的平行四边形定则可得：改变弹簧测力计 B 拉力的大小；减小重物 $M$ 的质量；改变弹簧测力计 B 拉力的方向，可以减小弹簧测力计 A 的拉力，或将 A 更换为较大的测力计。
}

\item
%% number: 11
%% paperName: 2011年江苏高考第 11 题 10 分
%% typeId: 4
%% type: 实验题
%% chapter: 电路
%% point: 电路实验
%% method: 无
%% score: 10
%% degree: 450
%% duplicateId: 0
%% body:
某同学利用如图所示的实验电路来测量电阻的阻值。

\onepicture[w=3.8cm]{figs/11b.png}

\begin{enumerate}
	\item
	将电阻箱接入 a、b 之间，闭合开关。适当调节滑动变阻器 $R'$ 后保持其阻值不变。改变电阻箱的阻值 $R$，得到一组电压表的示数 $U$ 与 $R$ 的数据如下表：

	\begin{center}
	\begin{tabular}{|c|c|c|c|c|c|c|}
	\hline
	电阻 $R/\UO$ & 5.0 & 10.0 & 15.0 & 25.0 & 35.0 & 45.0 \\
	\hline
	电压 $U/\UV$ & 1.00 & 1.50 & 1.80 & 2.14 & 2.32 & 2.45 \\
	\hline
	\end{tabular}
	\end{center}

	请根据实验数据作出 $U-R$ 关系图象。

	\onepicture[w=4.8cm, align=right]{figs/11a.png}
	\item
	用待测电阻 $R_{x}$ 替换电阻箱，读得电压表示数为 $2.00\UV$。利用（1）中测绘的 $U-R$ 图象可得 $R_{x}=$\tkanswer[1.2cm]{20} $\UO$。
	\item
	使用较长时间后，电池的电动势可认为不变，但内阻增大。若仍用本实验装置和（1）中测绘的 $U-R$ 图象测定某一电阻，则测定结果将\tkanswer[1.2cm]{偏小}（选填“偏大”或“偏小”）。现将一已知阻值为 $10\UO$ 的电阻换接在 a、b 之间，你应如何调节滑动变阻器，便仍可利用本实验装置和（1）中测绘的 $U-R$ 图象实现对待测电阻的准确测定？
\end{enumerate}

%% answer:
\jdanswer{
\begin{enumerate}
	\item
	根据实验数据描点作出 $U-R$ 关系图象，见图。
	\item
	$20$（读取 $19\sim21\UO$ 均可）
	\item
	偏小；把滑动变阻器 $R'$ 的阻值调小，使 $R'+r$ 保持不变。
\end{enumerate}
}
%% memo:
\memoanswer{
（1）根据实验数据描点作出 $U-R$ 关系图象，见图。

\onepicture[w=4.8cm]{figs/11_an.png}

（2）根据 $U-R$ 关系图象，电压表示数为 $2.00\UV$ 时，$R_{x}$ 的阻值为 $20\UO$，读取 $19\sim21\UO$ 可以。

（3）因为电压表示数 $U=IR=\frac{ER}{R+r+R'}=\frac{E}{1+\frac{r+R'}{R}}$，当电池的内阻 $r$ 增大时，同一个 $R$，则电压表读数将变小，按原来的 $U-R$ 图象，则电阻的测量值小于真实值，即偏小。要使电压表读数为 $1.50\UV$，因为电池内阻 $r$ 增大，应该把滑动变阻器阻值 $R'$ 调小，以至于使 $R'+r$ 不变。
}

\item
%% number: 12
%% paperName: 2011年江苏高考第 12 题 4 分
%% typeId: 8
%% type: 简答题
%% chapter: 原子和原子核
%% point: 质能方程
%% method: 无
%% score: 4
%% degree: 650
%% duplicateId: 0
%% body:
本题为选做题，包括 A、B、C 三小题，请选定其中两题，并在相应的答题区域内作答。若三题都做，则按 A、B 两题评分。

\textbf{A．}（选修模块 3-3）（12 分）

\begin{enumerate}
	\item
	如题图所示，一演示用的“永动机”转轮由 5 根轻杆和转轴构成，轻杆的末端装有形状记忆合金制成的叶片，轻推转轮后，进入热水的叶片因伸展而“划水”，推动转轮转动。离开热水后，叶片形状迅速恢复，转轮因此能较长时间转动。下列说法正确的是\xzanswer{D}

	\onepicture[w=4.4cm]{\includesvg{figs/12a}}

	\fourchoices[answer=D]
	{转轮依靠自身惯性转动，不需要消耗外界能量}
	{转轮转动所需能量来自形状记忆合金自身}
	{转动的叶片不断搅动热水，水温升高}
	{叶片在热水中吸收的热量一定大于在空气中释放的热量}
	\item
	如图所示，内壁光滑的气缸水平放置。一定质量的理想气体被密封在气缸内，外界大气压强为 $p_{0}$。现对气缸缓慢加热，气体吸收热量 $Q$ 后，体积由 $V_{1}$ 增大为 $V_{2}$。则在此过程中，气体分子平均动能\tkanswer[1.2cm]{增大}（选填“增大”、“不变”或“减小”），气体内能变化了\tkanswer[2.4cm]{$Q-p_{0}(V_{2}-V_{1})$}。

	\onepicture[align=right]{figs/12b.png}
	\item
	某同学在进行“用油膜法估测分子的大小”的实验前，查阅数据手册得知：油酸的摩尔质量 $M=0.283\Ukgmol$，密度 $\rho=0.895\times10^{3}\Ukgmc$。若 100 滴油酸的体积为 $1\UmL$，则 1 滴油酸所能形成的单分子油膜的面积约是多少？（取 $N_{A}=6.02\times10^{23}\Umol^{-1}$，球的体积 $V$ 与直径 $D$ 的关系为 $V=\frac{1}{6}\pi D^{3}$，结果保留一位有效数字）
\end{enumerate}

\textbf{B．}（选修模块 3-4）（12 分）

\begin{enumerate}
	\item
	如图所示，沿平直铁路线有间距相等的三座铁塔 A、B 和 C。假想有一列车沿 AC 方向以接近光速行驶，当铁塔 B 发出一个闪光，列车上的观测者测得 A、C 两铁塔被照亮的顺序是\xzanswer{C}

	\onepicture[w=5.4cm]{\includesvg{figs/12c}}

	\fourchoices[answer=C]
	{同时被照亮}
	{A 先被照亮}
	{C 先被照亮}
	{无法判断}
	\item
	一束光从空气射向折射率为 $\sqrt{3}$ 的某种介质，若反射光线与折射光线垂直，则入射角为\tkanswer[1.4cm]{$60^{\circ}$}。真空中的光速为 $c$，则光在该介质中的传播速度为\tkanswer[2.0cm]{$\frac{\sqrt{3}}{3}c$}。
	\item
	将一劲度系数为 $k$ 的轻质弹簧竖直悬挂，下端系上质量为 $m$ 的物块，将物块向下拉离平衡位置后松开，物块上下做简谐运动，其振动周期恰好等于以物块平衡时弹簧的伸长量为摆长的单摆周期。请由单摆周期公式推算出物块做简谐运动的周期 $T$。
\end{enumerate}

\textbf{C．}（选修模块 3-5）（12 分）

\begin{enumerate}
	\item
	下列描绘两种温度下黑体辐射强度与波长关系的图中，符合黑体辐射规律的是\xzanswer{A}

	\fourchoices[answer=A, ispicture=true, h=2.0cm]
	{figs/12da.png}
	{figs/12db.png}
	{figs/12dc.png}
	{figs/12dd.png}
	\item
	按照玻尔原子理论，氢原子中的电子离原子核越远，氢原子的能量\tkanswer[1.2cm]{越大}（选填“越大”或“越小”）。已知氢原子的基态能量为 $E_{1}$（$E_{1}<0$），电子质量为 $m$，基态氢原子中的电子吸收一频率为 $\nu$ 的光子被电离后，电子速度大小为\tkanswer[2.8cm]{$\sqrt{\frac{2(h\nu+E_{1})}{m}}$}（普朗克常量为 $h$）。
	\item
	有些核反应过程是吸收能量的。例如在 $\ce{X + ^{14}_{7}N -> ^{17}_{8}O + ^{1}_{1}H}$ 中，核反应吸收的能量 $Q=[(m_{O}+m_{H})-(m_{X}+m_{N})]c^{2}$，在该核反应中，X 表示什么粒子？X 粒子以动能 $E_{k}$ 轰击静止的 $^{14}_{7}N$，若 $E_{k}=Q$，则该核反应能否发生？请简要说明理由。
\end{enumerate}

%% answer:
\jdanswer{
\begin{enumerate}
	\item
	A：（1）D；（2）增大，$Q-p_{0}(V_{2}-V_{1})$；（3）$S=10\Umq$。
	\item
	B：（1）C；（2）$60^{\circ}$，$\frac{\sqrt{3}}{3}c$；（3）$T=2\pi\sqrt{\frac{m}{k}}$。
	\item
	C：（1）A；（2）越大，$\sqrt{\frac{2(h\nu+E_{1})}{m}}$；（3）$\ce{^{4}_{2}He}$；不能实现，因为不能同时满足能量守恒和动量守恒的要求。
\end{enumerate}
}
%% memo:
\memoanswer{
\textbf{A．}（1）轻推转轮后，叶片开始转动，由能量守恒定律可知，叶片在热水中吸收的热量在空气中释放和使叶片在热水中的膨胀做功，所以叶片在热水中吸收的热量一定大于在空气中释放的热量，D 正确。

（2）由于对气缸缓慢加热，温度升高，气体分子平均动能增大；根据热力学第一定律 $W+Q=\Delta E$，其中气体对外做功 $W=-p_{0}(V_{2}-V_{1})$。气体内能变化 $\Delta E=Q-p_{0}(V_{2}-V_{1})$。

（3）一个油酸分子的体积 $V=\frac{M}{\rho N_{A}}$，由球的体积与直径的关系得分子直径 $D=\sqrt{\frac{6M}{\pi\rho N_{A}}}$，最大面积 $S=\frac{1\times10^{-8}\Umc}{D}$，解得：$S=10\Umq$。

\textbf{B．}（1）当铁塔 B 发出一个闪光，同时到达 A、C 两铁塔被反射，但列车沿 AC 方向以接近光速行驶，经铁塔 A 反射的光相对列车的速度远小于经铁塔 C 反射的光相对列车的速度，经铁塔 C 反射的光先到达观测者，看到 C 先被照亮，C 正确。

（2）根据折射定律 $n=\frac{\sin i}{\sin r}$ 及 $i+r=90^{\circ}$，得 $\tan i=\sqrt{3}$，$i=60^{\circ}$；根据 $n=\frac{c}{v}$，得 $v=\frac{\sqrt{3}}{3}c$。

（3）单摆周期公式 $T=2\pi\sqrt{\frac{l}{g}}$，且 $kl=mg$，解得 $T=2\pi\sqrt{\frac{m}{k}}$。

\textbf{C．}（1）黑体辐射规律：随着温度的升高，一方面，各种波长的辐射强度都有增加，另一方面，辐射强度的极大值向波长较短的方向移动，A 图正确。

（2）根据 $E_{n}=\frac{-13.6\UeV}{n^{2}}$ 及 $r_{n}=n^{2}r_{0}$，越远，n 越大，$E_{n}$ 越大（注意 $E_{n}$ 为负值）。电子动能 $E_{k}=h\nu+E_{1}=\frac{1}{2}mv^{2}$（注意 $E_{1}$ 为负值），可解得：$v=\sqrt{\frac{2(h\nu+E_{1})}{m}}$。

（3）根据核反应的质量数和电荷数守恒，可得：X 粒子的质量数为 4，电荷数为 2，是 $\ce{^{4}_{2}He}$。
}

\item
%% number: 13
%% paperName: 2011年江苏高考第 13 题 15 分
%% typeId: 6
%% type: 计算题
%% chapter: 交流电
%% point: 变压器
%% method: 无
%% score: 15
%% degree: 450
%% duplicateId: 0
%% body:
如图为一理想变压器，ab 为原线圈，ce 为副线圈，d 为副线圈引出的一个接头。原线圈输入正弦式交变电压的 $u-t$ 图象如图所示。若只在 $ce$ 间接一只 $R_{ce}=400\UO$ 的电阻，或只在 de 间接一只 $R_{de}=225\UO$ 的电阻，两种情况下电阻消耗的功率均为 $80\UW$。

\begin{enumerate}
	\item
	请写出原线圈输入电压瞬时值 $u_{ab}$ 的表达式；
	\item
	求只在 ce 间接 $400\UO$ 的电阻时，原线圈中的电流 $I_{1}$；
	\item
	求 ce 和 de 间线圈的匝数比。
\end{enumerate}

\twopicture[showlabel=false, wA=3.6cm, wB=4.6cm, align=right]{figs/13a.png}{figs/13b.png}

%% answer:
\jdanswer{
\begin{enumerate}
	\item
	$u_{ab}=400\sin200\pi t$（$\UV$）
	\item
	$I_{1}=0.28\UA$（或 $\frac{\sqrt{2}}{5}\UA$）
	\item
	$\frac{4}{3}$
\end{enumerate}
}
%% memo:
\memoanswer{
（1）由题图知 $\omega=200\pi\Urads$，电压瞬时值 $u_{ab}=400\sin200\pi t$（$\UV$）。

（2）电压有效值 $U_{1}=200\sqrt{2}\UV$，理想变压器 $P_{1}=P_{2}$，原线圈中的电流 $I_{1}=\frac{P_{1}}{U_{1}}$，解得：$I_{1}\approx0.28\UA$（或 $\frac{\sqrt{2}}{5}\UA$）。

（3）设 ab 间匝数为 $n_{1}$，$\frac{U_{1}}{n_{1}}=\frac{U_{de}}{n_{de}}$，由题意知 $\frac{U_{ce}^{2}}{R_{ce}}=\frac{U_{de}^{2}}{R_{de}}$，解得：$\frac{n_{ce}}{n_{de}}=\sqrt{\frac{R_{ce}}{R_{de}}}$，代入数据得 $\frac{n_{ce}}{n_{de}}=\frac{4}{3}$。
}

\item
%% number: 14
%% paperName: 2011年江苏高考第 14 题 16 分
%% typeId: 6
%% type: 计算题
%% chapter: 曲线运动
%% point: 平抛运动
%% method: 无
%% score: 16
%% degree: 250
%% duplicateId: 0
%% body:
如图所示，长为 $L$、内壁光滑的直管与水平地面成 $30^{\circ}$ 角固定放置。将一质量为 $m$ 的小球固定在管底，用一轻质光滑细线将小球与质量为 $M=km$ 的小物块相连，小物块悬挂于管口。现将小球释放，一段时间后，小物块落地静止不动，小球继续向上运动，通过管口的转向装置后做平抛运动，小球在转向过程中速率不变。（重力加速度为 $g$）

\begin{enumerate}
	\item
	求小物块下落过程中的加速度大小；
	\item
	求小球从管口抛出时的速度大小；
	\item
	试证明小球平抛运动的水平位移总小于 $\frac{\sqrt{2}}{2}L$。
\end{enumerate}

\onepicture[w=6.0cm, align=right]{figs/14.png}

%% answer:
\jdanswer{
\begin{enumerate}
	\item
	$a=\frac{2k-1}{2(k+1)}g$
	\item
	$v_{0}=\sqrt{\frac{k-2}{2(k+1)}gL}$，（$k>2$）
	\item
	平抛运动 $x=v_{0}t$，$L\sin30^{\circ}=\frac{1}{2}gt^{2}$，解得 $x=\sqrt{\frac{k-2}{2(k+1)}}L$；因为 $\frac{k-2}{k+1}<1$，所以 $x<\frac{\sqrt{2}}{2}L$，得证。
\end{enumerate}
}
%% memo:
\memoanswer{
（1）设细线中的张力为 $T$，根据牛顿第二定律

$Mg-T=Ma$

$T-mg\sin30^{\circ}=ma$

且 $M=km$

解得：$a=\frac{2k-1}{2(k+1)}g$。

（2）设 M 落地时的速度大小为 $v$，m 射出管口时速度大小为 $v_{0}$，M 落地后 m 的加速度为 $a_{0}$。根据牛顿第二定律 $-mg\sin30^{\circ}=ma_{0}$

匀变速直线运动 $v^{2}=2aL\sin30^{\circ}$

$v_{0}^{2}-v^{2}=2a_{0}L(1-\sin30^{\circ})$

解得：$v_{0}=\sqrt{\frac{k-2}{2(k+1)}gL}$（$k>2$）。

（3）平抛运动 $x=v_{0}t$，$L\sin30^{\circ}=\frac{1}{2}gt^{2}$，解得 $x=L\sqrt{\frac{k-2}{2(k+1)}}$。

因为 $\frac{k-2}{k+1}<1$，所以 $x<\frac{\sqrt{2}}{2}L$，得证。
}

\item
%% number: 15
%% paperName: 2011年江苏高考第 15 题 16 分
%% typeId: 6
%% type: 计算题
%% chapter: 磁场
%% point: 带电粒子在组合场中的运动
%% method: 无
%% score: 16
%% degree: 250
%% duplicateId: 0
%% body:
某种加速器的理想模型如图 \ref{2011江苏15a} 所示：两块相距很近的平行小极板中间各开有一小孔 a、b，两极板间电压 $u_{ab}$ 的变化图象如图 \ref{2011江苏15b} 所示，电压的最大值为 $U_{0}$、周期为 $T_{0}$，在两极板外有垂直纸面向里的匀强磁场。若将一质量为 $m_{0}$、电荷量为 $q$ 的带正电的粒子从板内 a 孔处静止释放，经电场加速后进入磁场，在磁场中运动时间 $T_{0}$ 后恰能再次从 a 孔进入电场加速。现该粒子的质量增加了 $\frac{1}{100}m_{0}$。（粒子在两极板间的运动时间不计，两极板外无电场，不考虑粒子所受的重力）

\begin{enumerate}
	\item
	若在 $t=0$ 时刻将该粒子从板内 $a$ 孔处静止释放，求其第二次加速后从 b 孔射出时的动能；
	\item
	现在利用一根长为 $L$ 的磁屏蔽管（磁屏蔽管置于磁场中时管内无磁场，忽略其对管外磁场的影响），使图 \ref{2011江苏15a} 中实线轨迹（圆心为 O）上运动的粒子从 a 孔正下方相距 $L$ 处的 c 孔水平射出，请在答题卡图上的相应位置处画出磁屏蔽管；
	\item
	若将电压 $u_{ab}$ 的频率提高为原来的 2 倍，该粒子应何时由板内 a 孔处静止开始加速，才能经多次加速后获得最大动能？最大动能是多少？
\end{enumerate}

\twopicture[numstyle=arabic, labelA=2011江苏15a, labelB=2011江苏15b, wA=4.6cm, wB=5.8cm, align=right]{figs/15a.png}{figs/15b.png}

%% answer:
\jdanswer{
\begin{enumerate}
	\item
	$E_{k2}=\frac{49}{25}qU_{0}$
	\item
	磁屏蔽管的位置如图所示。
	\item
	$E_{km}=\frac{313}{25}qU_{0}$
\end{enumerate}
}
%% memo:
\memoanswer{
（1）质量为 $m_{0}$ 的粒子在磁场中作匀速圆周运动 $Bqv=\frac{mv^{2}}{r}$，$T_{0}=\frac{2\pi r}{v}$

则 $T_{0}=\frac{2\pi m_{0}}{qB}$

当粒子的质量增加了 $\frac{1}{100}m_{0}$，其周期增加 $\Delta T=\frac{1}{100}T_{0}$

根据题图可知，粒子第一次的加速电压 $u_{1}=U_{0}$

粒子第二次的加速电压 $u_{2}=\frac{24}{25}U_{0}$

粒子射出时的动能 $E_{k2}=qu_{1}+qu_{2}$

解得 $E_{k2}=\frac{49}{25}qU_{0}$。

（2）磁屏蔽管的位置如图所示。

\onepicture[w=2.6cm, align=right]{figs/15_an.png}

（3）在 $u_{ab}>0$ 时，粒子被加速，则最多连续被加速的次数

$N=\frac{\frac{T_{0}}{4}}{\Delta T}$，得 $N=25$

分析可得，粒子在连续被加速的次数最多，且 $u=U_{0}$ 时也被加速的情况时，最终获得的动能最大。

粒子由静止开始被加速的时刻 $t=\left(\frac{n}{2}+\frac{19}{50}\right)T_{0}$（$n=0,1,2,\ldots$）

最大动能 $E_{km}=2\times\left(\frac{1}{25}+\frac{3}{25}+\cdots\frac{23}{25}\right)qU_{0}+qU_{0}$ 解得 $E_{km}=\frac{313}{25}qU_{0}$。
}

\end{enumerate}

\end{document}
